% notes-l-structures.tex
% Volume I — Model Theory
% Topic: L-Structures (Languages and their interpretations)
% Status: Active

% ──────────────────────────────────────────────────────────────
%  TOOLKIT — L-Structures Quick Reference
% ──────────────────────────────────────────────────────────────

\begin{tcolorbox}[colback=gray!6, colframe=gray!40, arc=2pt,
  left=6pt, right=6pt, top=4pt, bottom=4pt,
  title={\small\textbf{L-Structures Toolkit — Quick Reference}},
  fonttitle=\small\bfseries]
\small
\begin{tabular}{l l l}
\toprule
\textbf{Concept} & \textbf{Symbol / Notation} & \textbf{Detail} \\
\midrule
First-order language    & $\mathcal{L}$                          & \hyperref[def:first-order-language]{↓} \\
L-structure             & $\mathcal{M} = (M, \mathcal{L}^{\mathcal{M}})$ & \hyperref[def:l-structure]{↓} \\
Domain / universe       & $|\mathcal{M}|$ or $M$                 & \hyperref[def:l-structure]{↓} \\
Interpretation of $f$   & $f^{\mathcal{M}} : M^n \to M$          & \hyperref[def:l-structure]{↓} \\
Interpretation of $R$   & $R^{\mathcal{M}} \subseteq M^n$        & \hyperref[def:l-structure]{↓} \\
Interpretation of $c$   & $c^{\mathcal{M}} \in M$                & \hyperref[def:l-structure]{↓} \\
Variable assignment     & $s : \mathrm{Var} \to M$               & \hyperref[def:variable-assignment]{↓} \\
Term evaluation         & $s(t) \in M$                           & \hyperref[def:term-evaluation]{↓} \\
Satisfaction            & $\mathcal{M} \models \varphi[s]$       & \hyperref[def:satisfaction]{↓} \\
\bottomrule
\end{tabular}
\end{tcolorbox}

\vspace{0.5em}

% ──────────────────────────────────────────────────────────────
%  §1  FIRST-ORDER LANGUAGES
% ──────────────────────────────────────────────────────────────

\section{First-Order Languages}

\begin{tcolorbox}[colback=propbox, colframe=propborder, arc=2pt,
  left=6pt, right=6pt, top=4pt, bottom=4pt,
  title={\small\textbf{Definition (First-Order Language $\mathcal{L}$)}},
  fonttitle=\small\bfseries]
\label{def:first-order-language}
A \textbf{first-order language} $\mathcal{L}$ consists of three (possibly
empty) sets of non-logical symbols:
\begin{itemize}
  \item a set $\mathcal{R}$ of \textbf{relation symbols} (also called
        \emph{predicate symbols}), each assigned a fixed arity $n \geq 1$;
  \item a set $\mathcal{F}$ of \textbf{function symbols}, each assigned a
        fixed \textbf{arity} $n \geq 0$ (0-ary function symbols are
        \emph{constant symbols});
  \item a set $\mathcal{C}$ of \textbf{constant symbols} (equivalently,
        0-ary function symbols — listed separately for emphasis).
\end{itemize}
Together with the logical symbols common to all first-order languages
(variables, connectives $\lnot, \land, \lor, \to, \leftrightarrow$,
quantifiers $\forall, \exists$, equality $=$, parentheses), the triple
$(\mathcal{F}, \mathcal{R}, \mathcal{C})$ determines $\mathcal{L}$.
\end{tcolorbox}

\begin{remark*}[Fully quantified form]
$\mathcal{L} = (\mathcal{F}, \mathcal{R}, \mathcal{C})$ where
$\mathrm{ar}(f) \in \mathbb{N}$ for each $f \in \mathcal{F}$ and
$\mathrm{ar}(R) \in \mathbb{N}_{\geq 1}$ for each $R \in \mathcal{R}$.
\end{remark*}

\begin{remark*}[Intuition]
A language is a \emph{vocabulary}: it names the operations, relations, and
distinguished elements that can appear in sentences, without yet saying what
those names mean. The meaning is supplied by a structure.
\end{remark*}

\begin{remark*}[Source note]
Some sources (e.g.\ Marker \S1.1, Hodges) absorb constant symbols into
0-ary function symbols. Others (e.g.\ Chang--Keisler) list them separately.
The content is equivalent; only the bookkeeping differs.
\end{remark*}

\begin{remark*}[Consequence]
The set of $\mathcal{L}$-terms and the set of $\mathcal{L}$-formulas are
both defined inductively from $\mathcal{L}$; see
\hyperref[def:term-evaluation]{term evaluation} below.
\end{remark*}

% ── Examples of languages ─────────────────────────────────────

\begin{example}[Standard examples of first-order languages]
\begin{enumerate}
  \item \textbf{Language of ordered fields} $\mathcal{L}_{\mathrm{of}}$:
        function symbols $\{+, \cdot, -\}$ of arities $2, 2, 1$;
        relation symbol $\{<\}$ of arity $2$;
        constant symbols $\{0, 1\}$.
        The real field $(\mathbb{R}, +, \cdot, -, <, 0, 1)$ is an
        $\mathcal{L}_{\mathrm{of}}$-structure.

  \item \textbf{Language of graphs} $\mathcal{L}_{\mathrm{gr}}$:
        no function symbols, no constants;
        one binary relation symbol $\{E\}$ (for ``edge'').

  \item \textbf{Language of groups} $\mathcal{L}_{\mathrm{gp}}$:
        function symbols $\{\cdot, {}^{-1}\}$ of arities $2, 1$;
        constant symbol $\{e\}$.

  \item \textbf{Pure equality language} $\mathcal{L}_{=}$:
        no non-logical symbols at all. Structures are arbitrary
        non-empty sets.
\end{enumerate}
\end{example}

% ──────────────────────────────────────────────────────────────
%  §2  L-STRUCTURES
% ──────────────────────────────────────────────────────────────

\subsection{L-Structures}

\begin{tcolorbox}[colback=propbox, colframe=propborder, arc=2pt,
  left=6pt, right=6pt, top=4pt, bottom=4pt,
  title={\small\textbf{Definition ($\mathcal{L}$-Structure)}},
  fonttitle=\small\bfseries]
\label{def:l-structure}
Let $\mathcal{L} = (\mathcal{F}, \mathcal{R}, \mathcal{C})$ be a
first-order language.
An \textbf{$\mathcal{L}$-structure} $\mathcal{M}$ consists of:
\begin{enumerate}
  \item a non-empty set $M$ called the \textbf{domain} (or
        \textbf{universe}) of $\mathcal{M}$, written $|\mathcal{M}|$;
  \item for each $n$-ary relation symbol $R \in \mathcal{R}$,
        a relation $R^{\mathcal{M}} \subseteq M^n$;
  \item for each $n$-ary function symbol $f \in \mathcal{F}$,
        a function $f^{\mathcal{M}} : M^n \to M$;
  \item for each constant symbol $c \in \mathcal{C}$,
        a distinguished element $c^{\mathcal{M}} \in M$.
\end{enumerate}
The superscript notation $f^{\mathcal{M}}, R^{\mathcal{M}}, c^{\mathcal{M}}$
denotes the \textbf{interpretation} of each symbol in $\mathcal{M}$.
\end{tcolorbox}

\begin{remark*}[Fully quantified form]
$\mathcal{M}$ is an $\mathcal{L}$-structure iff
$M \neq \emptyset$,
and $\forall f \in \mathcal{F}_{(n)},\ f^{\mathcal{M}} : M^n \to M$,
and $\forall R \in \mathcal{R}_{(n)},\ R^{\mathcal{M}} \subseteq M^n$,
and $\forall c \in \mathcal{C},\ c^{\mathcal{M}} \in M$.
\end{remark*}

\begin{remark*}[Intuition]
A structure is an $\mathcal{L}$-vocabulary \emph{made concrete}: each
symbol is assigned an actual mathematical object of the correct type.
The domain provides the raw material; the interpretations give the
symbols their meaning inside that domain.
\end{remark*}

\begin{remark*}[Common error]
The domain $M$ must be \emph{non-empty}. First-order logic with an empty
domain is pathological: $\forall x\, \varphi$ would be vacuously true
and $\exists x\, \varphi$ vacuously false for every $\varphi$, breaking
the duality of the quantifiers.
\end{remark*}

\begin{remark*}[Source note]
Some sources write $\mathfrak{M}$ (Fraktur) for structures and reserve
$M$ for the domain; others use $\mathcal{M}$ throughout and write
$|\mathcal{M}|$ for the domain. The latter convention is used here.
Marker uses $\mathcal{M}$; Enderton uses $\mathfrak{A}$.
\end{remark*}

% ── Examples of L-structures ──────────────────────────────────

\begin{example}[L-structures for familiar mathematical objects]
\begin{enumerate}
  \item $(\mathbb{R}, +^{\mathbb{R}}, \cdot^{\mathbb{R}}, -^{\mathbb{R}},
         <^{\mathbb{R}}, 0^{\mathbb{R}}, 1^{\mathbb{R}})$
        is an $\mathcal{L}_{\mathrm{of}}$-structure with domain
        $M = \mathbb{R}$.

  \item $(\mathbb{Q}, +^{\mathbb{Q}}, \cdot^{\mathbb{Q}}, -^{\mathbb{Q}},
         <^{\mathbb{Q}}, 0^{\mathbb{Q}}, 1^{\mathbb{Q}})$
        is another $\mathcal{L}_{\mathrm{of}}$-structure with domain
        $M = \mathbb{Q}$.

  \item The complete graph $K_3$ on $\{1,2,3\}$ is an
        $\mathcal{L}_{\mathrm{gr}}$-structure with
        $E^{K_3} = \{(1,2),(2,1),(1,3),(3,1),(2,3),(3,2)\}$.
\end{enumerate}
Two different structures can share the same language — the language only
fixes the \emph{type} of the interpretations, not their specific content.
\end{example}

% ──────────────────────────────────────────────────────────────
%  §3  VARIABLE ASSIGNMENTS
% ──────────────────────────────────────────────────────────────

\subsection{Variable Assignments}

\begin{tcolorbox}[colback=propbox, colframe=propborder, arc=2pt,
  left=6pt, right=6pt, top=4pt, bottom=4pt,
  title={\small\textbf{Definition (Variable Assignment)}},
  fonttitle=\small\bfseries]
\label{def:variable-assignment}
Let $\mathcal{M}$ be an $\mathcal{L}$-structure with domain $M$, and let
$\mathrm{Var} = \{v_1, v_2, v_3, \ldots\}$ be a fixed countable set of
first-order variables.

A \textbf{variable assignment} (or \textbf{environment}) for $\mathcal{M}$
is a function $s : \mathrm{Var} \to M$.

For a variable assignment $s$, a variable $v_i$, and an element $a \in M$,
write $s[v_i \mapsto a]$ for the assignment that agrees with $s$ on all
variables except $v_i$, where it returns $a$:
\[
  s[v_i \mapsto a](v_j) =
  \begin{cases}
    a      & \text{if } j = i, \\
    s(v_j) & \text{if } j \neq i.
  \end{cases}
\]
\end{tcolorbox}

\begin{remark*}[Fully quantified form]
$s : \mathrm{Var} \to M$, and
$s[v_i \mapsto a](v_j) = a$ if $j = i$, else $s(v_j)$.
\end{remark*}

\begin{remark*}[Intuition]
A variable assignment gives temporary values to free variables before
evaluating a formula. The $s[v_i \mapsto a]$ notation is the standard
tool for handling quantifier rules: to check $\exists v_i\,\varphi$,
test whether $\mathcal{M} \models \varphi[s[v_i \mapsto a]]$ for some
$a \in M$.
\end{remark*}

% ──────────────────────────────────────────────────────────────
%  §4  TERM EVALUATION
% ──────────────────────────────────────────────────────────────

\subsection{Term Evaluation}

\begin{tcolorbox}[colback=propbox, colframe=propborder, arc=2pt,
  left=6pt, right=6pt, top=4pt, bottom=4pt,
  title={\small\textbf{Definition (Term Evaluation)}},
  fonttitle=\small\bfseries]
\label{def:term-evaluation}
Let $\mathcal{M}$ be an $\mathcal{L}$-structure and $s : \mathrm{Var} \to M$
a variable assignment.
The \textbf{evaluation} $s(t) \in M$ of an $\mathcal{L}$-term $t$ under
$s$ is defined inductively:
\begin{enumerate}
  \item If $t = v_i$ is a variable, then $s(t) = s(v_i)$.
  \item If $t = c$ is a constant symbol, then $s(t) = c^{\mathcal{M}}$.
  \item If $t = f(t_1, \ldots, t_n)$ for an $n$-ary function symbol $f$,
        then $s(t) = f^{\mathcal{M}}(s(t_1), \ldots, s(t_n))$.
\end{enumerate}
\end{tcolorbox}

\begin{remark*}[Fully quantified form]
$s : \mathrm{Var} \to M$, and $s(t) \in M$ for every $\mathcal{L}$-term $t$.
Formally: $s(\cdot)$ is the unique function on terms satisfying clauses (1)--(3).
\end{remark*}

\begin{remark*}[Intuition]
Term evaluation is the ``computation'' step: given concrete values for
all variables, every term reduces to an element of $M$. The inductive
definition mirrors the inductive definition of terms themselves.
\end{remark*}

% ──────────────────────────────────────────────────────────────
%  §5  SATISFACTION  [STUB]
% ──────────────────────────────────────────────────────────────

\subsection{Satisfaction \textnormal{[Stub]}}

\begin{tcolorbox}[colback=propbox, colframe=propborder, arc=2pt,
  left=6pt, right=6pt, top=4pt, bottom=4pt,
  title={\small\textbf{Definition (Satisfaction, $\mathcal{M} \models \varphi[s]$) — Stub}},
  fonttitle=\small\bfseries]
\label{def:satisfaction}
Let $\mathcal{M}$ be an $\mathcal{L}$-structure, $s$ a variable assignment,
and $\varphi$ an $\mathcal{L}$-formula.
The \textbf{satisfaction relation} $\mathcal{M} \models \varphi[s]$
(read: ``$\mathcal{M}$ satisfies $\varphi$ under $s$'') is defined
inductively on the structure of $\varphi$:

\medskip
\textit{[To be filled in — clauses for atomic formulas, negation,
conjunction, disjunction, implication, existential quantifier,
universal quantifier.]}
\end{tcolorbox}

\begin{remark*}[Fully quantified form]
\textit{[Stub — to be completed.]}
\end{remark*}

\begin{remark*}[Intuition]
Satisfaction is the semantic notion of truth: $\mathcal{M} \models \varphi[s]$
means $\varphi$ is true in $\mathcal{M}$ when each free variable $v_i$ takes
the value $s(v_i)$.
The definition proceeds by induction on formula complexity,
grounding out on atomic formulas evaluated via term evaluation.
\end{remark*}

% ──────────────────────────────────────────────────────────────
%  §6  ELEMENTARY EQUIVALENCE AND SUBSTRUCTURES  [STUB]
% ──────────────────────────────────────────────────────────────

\subsection{Elementary Equivalence and Substructures \textnormal{[Stub]}}

\begin{tcolorbox}[colback=gray!6, colframe=gray!40, arc=2pt,
  left=6pt, right=6pt, top=4pt, bottom=4pt,
  title={\small\textbf{Status: Planned}},
  fonttitle=\small\bfseries]
The following topics are planned for this section:
\begin{enumerate}
  \item Substructures ($\mathcal{N} \subseteq \mathcal{M}$).
  \item Elementary substructures ($\mathcal{N} \prec \mathcal{M}$) and
        the Tarski--Vaught criterion.
  \item Elementary equivalence ($\mathcal{M} \equiv \mathcal{N}$).
  \item Isomorphisms of $\mathcal{L}$-structures.
\end{enumerate}
\end{tcolorbox}
